%% CPP 2027 Submission — 0pat Exercise XXXII: Imaginary-Axis Coefficient Type and the Origin
\documentclass[sigplan,10pt,anonymous,review]{acmart}
\settopmatter{printfolios=true,printacmref=false}

\AtBeginDocument{%
  \providecommand\BibTeX{{\normalfont B\scshape ib\TeX}}}

\settopmatter{printacmref=false}
\renewcommand\footnotetextcopyrightpermission[1]{}
\pagestyle{plain}

\usepackage{microtype}
\usepackage{booktabs}
\usepackage{amsmath}
% XeTeX (tectonic): acmart loads newtxmath/libertine which already defines \Bbbk;
% reset it so amssymb's definition is accepted.
\let\Bbbk\relax
\usepackage{amssymb}

\begin{document}

\title{Imaginary-Axis Coefficient Type and the Origin\\
(Exercise XXXII)}

\author{Anonymous}
\affiliation{%
  \institution{Anonymous}
}
\email{anonymous@example.org}

\begin{abstract}
We compute the origin of the split-complex Euler circle frame with the
coefficient type of the imaginary axis made explicit. The imaginary
axis represents the imaginary part of the complex numbers; its
coefficients may be taken as natural numbers (discrete axis
$\{n \cdot i : n \in \mathbb{N}\}$) or as real numbers (continuous axis
$\{t \cdot i : t \in \mathbb{R}\}$). Four machine-checked theorems:
under natural coefficients, the real axis meets the imaginary axis at
$\{(0,0)\}$ (solving $x = n \cdot i$ forces $x = 0$ and $n = 0$);
under real coefficients, the intersection is again $\{(0,0)\}$
(solving $x = t \cdot i$ forces $x = 0$ and $t = 0$); the transposition
maps the real axis onto the imaginary axis --- the transposed real axis
is the imaginary axis with real (continuous) coefficients; and the
transposition preserves natural coefficients (the natural imaginary
axis maps to the natural real axis). The origin is $(0,0)$ under both
coefficient conventions. Zero errors, zero \texttt{sorry}.
\end{abstract}

\keywords{Lean, formalization, split-complex, coefficient type, origin, transposition, 0pat}

\maketitle

\section{The computation}

The orthogonal Euler circle frame (Exercises XXVII--XXXI) lives in the
split-complex plane $a + bj$, $j^2 = +1$. The imaginary axis represents
the imaginary part of the complex numbers. This exercise makes the
coefficient type of the imaginary axis explicit and computes the origin
under each convention.

\section{The two computations}

\paragraph{Under natural coefficients (origin\_natural\_coefficients).}
The imaginary axis is the discrete set $\{n \cdot i : n \in
\mathbb{N}\}$; the real axis is the continuous real line. A point on
both axes has the form $x$ (real axis) and $n \cdot i$ (natural
imaginary axis); solving $x = n \cdot i$ forces $x = 0$ and $n = 0$.
The intersection is the single point $(0,0)$. Under natural
coefficients, the origin is $(0,0)$.

\paragraph{Under imaginary-part (real) coefficients
(origin\_imaginary\_part).} The imaginary axis is the continuous set
$\{t \cdot i : t \in \mathbb{R}\}$. Solving $x = t \cdot i$ forces
$x = 0$ and $t = 0$. The intersection is again the single point
$(0,0)$. Under real coefficients, the origin is $(0,0)$.

The two computations agree: the origin of the frame is $(0,0)$ whether
the imaginary axis is discrete (natural coefficients) or continuous
(real coefficients). The natural coefficient $n = 0$ and the real
coefficient $t = 0$ are the unique solutions of their equations.

\section{The transposition}

\paragraph{The real axis transposes to the imaginary axis
(transpose\_realAxis\_eq\_imaginaryAxis).} The transposition
$(a,b) \mapsto (b,a)$ maps the real axis $\{(x, 0)\}$ onto the
imaginary axis $\{(0, x)\}$: the transposed real axis is the imaginary
axis, with real (continuous) coefficients. If the coefficients of the
imaginary axis are locked to natural numbers, the transposed real axis
is the imaginary axis computed by imaginary part (continuous) --- the
transposition itself does not discretize.

\paragraph{Natural coefficients are preserved by transposition
(transpose\_natural\_imaginary\_is\_natural\_real).} The natural
imaginary axis $\{n \cdot i\}$ maps under transposition to the natural
real axis $\{(n, 0)\}$. Discrete stays discrete: transposition changes
the axis, not the coefficient type.

\section{Honest boundary and position}

The content is elementary algebra; the value is the explicit separation
of the two coefficient conventions and the proof that they agree on the
origin. The result concerns the frame; no claim is made about the
Riemann zeta function. The analysis of $\zeta$ remains the open span of
the bridge.

\section{Formalization}

The proof is a single Lean file (\texttt{proof.lean}, 175 lines),
importing only \texttt{Mathlib.Data.Real.Basic} and basic tactics. The
split-complex structure (30 lines, as in Exercises XXVII--XXXI) is
self-contained. Four theorems, compiled with Lean 4.32.2 / mathlib
v4.32.2, zero errors, zero \texttt{sorry}.

\begin{center}
\begin{tabular}{lll}
\toprule
Theorem & Statement & Proof core \\
\midrule
\texttt{origin\_natural\_coefficients} & axes meet at $\{(0,0)\}$ ($n \in \mathbb{N}$) & ext; cast \\
\texttt{origin\_imaginary\_part} & axes meet at $\{(0,0)\}$ ($t \in \mathbb{R}$) & ext; simp \\
\texttt{transpose\_realAxis\_eq\_imaginaryAxis} & $T(\text{real axis}) = \text{imaginary axis}$ & ext; simp \\
\texttt{transpose\_natural\_...} & natural coefficients preserved & ext; rcases \\
\bottomrule
\end{tabular}
\end{center}

\section{Conclusion}

With the coefficient type of the imaginary axis made explicit, the
origin of the Euler circle frame is computed twice --- under natural
coefficients and under imaginary-part (real) coefficients --- and both
computations give $(0,0)$. The transposition maps the real axis onto
the imaginary axis with continuous coefficients and preserves natural
coefficients. The frame is fully determined; the analysis of $\zeta$
remains the open span of the bridge.

\bibliographystyle{ACM-Reference-Format}
\bibliography{refs}

\end{document}
